% !TeX root = ../papers/nc-p3.tex
\section{階数二平面の四行一次表示と三重点からの射影}
\label{sec:rank-two-plane-bordered-quartic}
\begin{theorem}\label{thm:rank-two-plane-bordered-quartic}
被約な四辺有限の標準有限Frobenius窓の一般可逆な全階数反対根 $N$ について、
階数二平面 $H=(h=0)$ が階数一曲線を含むとする。平面上の固定像方向による商での三行切断を $N_0$ とする。
四次式 $\det N_0/h^2$ は恒等零か、重複既約因子を持つ。
\end{theorem}
\begin{proof}[一次表示と曲線上の二方向]
命題\ref{prop:rank-two-plane-conic-trace}の固定方向を最後の座標に取り、一次列の平面上の像を補空間とする。
平面上では $N$ は三行の平方 $ww^t$ と一行の非零二次式の直和で、$w$ は平面の独立な三座標である。
従って $H$ 上ではその二次式の零集合だけが階数一で、直線上三印、平面二次上六印の上界から
平面外の印数は $n\ge14$ である。$w$ を一次に延長して
\[
 N=\begin{pmatrix}N_0&h\ell\\ h\ell^t&F\end{pmatrix},
 \qquad N_0=ww^t+hC,\qquad
 L=\begin{pmatrix}C&w\\ w^t&-h\end{pmatrix}
\]
と書く。$\ell$ の三成分は独立、最後の列 $(w,-h)$ は四独立なので $L$ は線型な四空間で各点で非零である。
$\det L=-\det N_0/h^2$ であり、平面外の印で $\operatorname{rank} L=1+\operatorname{rank}N_0\le2$。

$L$ がどこでも階数一にならないとき、$N_0$ は全点で非零である（$H$ 上は $w^2$）。
$L$ が階数二以下となる既約曲線 $J$ が $H$ に含まれず次数 $e$ であるとする。
$J$ 上の $N_0$ の平方方向 $b$ は正規化上の射であり、その超平面束の次数は $e$ である。
実際二倍が $N_0$ の二次座標の束で基点はない。
印では $b\wedge\ell=0$ となる。これが恒等零でなければ次数 $2e$ の非零切断により
平面外の印は $2e$ 以下である。恒等零なら $\ell$ は正規化上で共通零点を持てない。
基点分を除いたその射は $b$ に等しく、既に次数 $e$ を要するからである。
従って $N_0=c\ell\ell^t$ が $J$ 上で成り立ち、$c$ は極も零もない定数である。
平面外では $N$ の階数一条件は $cF=h^2$ となる。この式が恒等的なら $J$ 自身が階数一曲線である。
この場合も $J$ は $H$ 上の既存の曲線とともに階数一になるが、ここでその可能性を確かめる。
$H$ 上に被約次数二があれば総次数上限に反する。二次式が重複直線なら、追加曲線は
その直線と交わる直線に限られる。この場合は $n\ge17$ で、追加のその一本だけで
$2e$ の代わりに三印を許せば十分である。
他の全ての場合に二次式の切断により上界 $2e$ を得る。
\end{proof}
\begin{proof}[整四次式]
まず $\det L$ が既約かつ被約であると仮定する。階数二以下はその特異軌跡に入り、
曲線部分の被約総次数 $d$ は三以下である。一般平面切断は整四次曲線で、各特異点が
算術種数三を少なくとも一つ使うからである。曲線部分の外の階数二以下の点は $10-d$ 以下である。
実際、対称四行の階数二以下の多様体の次数は十である。二因子の積を取る写像は
$\mathbb P^3\times\mathbb P^3$ から次数二で、$(1,1)^6=20$ を使えばよい。
$\mathbb P(L)$ を含む一般の四平面はこの多様体と固有に次数十の曲線で交わる。
これは $\mathbb P(L)$ の外での一般超平面切断と、その内側の次元上限一による。
$\mathbb P(L)$ 内に含まれる全成分を除くと次数は高々 $10-d$ で、残りをその超平面で切ればよい。
孤立点はこの残りに入る（線型切断の全成分が次元一下限を持つ）。
従って階数一点がなければ先の曲線評価と合わせて $n\le10+d\le13$。
重複直線の場合の追加一印も $n\ge17$ には届かない。

階数一点 $p$ がある場合を扱う。$p$ と階数二曲線の各点を結ぶ直線は和の階数が三以下なので
四次曲面に入る。次数三以下の非直線曲線との結合は平面か次数三以下の錐面であり不可能である。
直線成分も $p$ を通らなければ平面を張る。曲線がなければ次数十で終わるので、
残る $d$ 本は全て $p$ を通る直線である。
$p=(1:0:0:0)$ とし、合同変換で $L(p)$ を一座標の平方に取る。$p$ からの方向 $z\in\mathbb P^2$ を使い
\[
   L=\begin{pmatrix}t+c(z)&m(z)^t\\m(z)&B(z)\end{pmatrix},
   \qquad \det L=(t+c)\det B-m^t\operatorname{adj}(B)m
\]
とする。三次式 $\det B$ が恒等零の場合は次段落で別に扱う。
$\det B\ne0$ の場合は既約性より三次式と残余四次式に共通因子はない。
$B$ が階数二の方向では単元二行を消去すると、残余四次式の三次曲線上での制限は
局所的に単元と一つの式の平方の積である。従って階数二の入力に対応する方向の交数は二以上であり、
消去後の対角の一次式が $t$ を一意に決める。$B$ の階数が一以下で階数二以下の入力を持つ方向では、
$B$ が一なら $m$ はその像に入る必要があり、階数零なら任意の $m$ でよい。
いずれも母線全体が階数二以下となる。この方向では三次式が特異で交数は二以上である。
従って全方向数は $3\cdot4/2=6$ 以下、うち $d$ 方向が直線成分である。
印数は $6-d+3d\le12$ となる（$p$ は既に母線上である）。

$\det B=0$ が恒等的なとき、$B$ の一般階数は二でなければ四次式は零である。
三変数環で二次余因子を $a gg^t$ と原始に書くと四次式は $-a(m^tg)^2$ となり、既約四次ではない。
以上で整四次を除いた。
\end{proof}
\begin{proof}[可約かつ被約な四次式]
$\det L$ の各平面因子 $P$ には平面外印（ここでは常に $H$ の外を意味する）が六つ以下である。
実際重複因子でないので、$P$ 上では一般階数三であり、三次余因子を $a gg^t$ と原始に書ける。
$g$ が非定数なら一次で、$a$ は一次、$g$ の基点は高々一点である。
階数二以下はその直線と一点で四印以下である。
核が一定の $z$ なら全ての像は $z^\perp$ に入る。最後の座標ベクトルを $e$ として
$z^\perp\ne e^\perp$ である。実際 $L(x)e$ は独立四座標なので一平面を消さない。
平面外印の出力を $(y',y_4)$ とすれば $(y',0)$ は $L$ の像に入る。
これは最後の行列要素 $-h$ の消去により、その列空間に $N_0$ の像を最後の座標零で持ち上げられるからである。
従って $P$ の該当印の出力は $z^\perp\cap e^\perp$ の二空間と最後の出力座標が張る平面 $P'$ に入る。
先の素な二直線の証明と同様に、平方座標同型が $P,P'$ の二次の包を同定すると固定列の平面因子が生じる。
よって対応印は $P$ 内の非零二次式で消え、六印以下である。

分割 $(3,1)$ では平面外の特異点は整三次の特異点で、孤立なら四点以下、特異直線ならその上の三印以下
（三重頂点を持つ錐の場合も特異母線は高々一本）。従って平面因子と合わせても十四に足りない。
$(2,1,1)$ では両平面外に整二次曲面の頂点一つまで、合計十三以下である。
$(1,1,1,1)$ では各特異点が二平面以上に入り、六印ずつの上界で十二以下となる。

$(2,2)$ で両因子の交叉を $Y$ とする。その外の特異点は高々二つの頂点である。
$L$ に階数一点があればその点では四次式が三重以上なので、一方の頂点が他方の曲面に入る。
よって $Y$ の外は一つ以下である。また $Y$ には十二印までしか載らない。
二つの互いに素な二次式の一次式倍は八独立な三次式を与え、三次評価は同型だからである。
これで階数一点のある場合を除いた。階数一点がなければ階数二曲線の被約総次数 $d$ は四以下で、
前段の $10-d$ と $2d$ の評価を使える。$Y$ が全て階数二曲線で覆われない場合でも、
$d\le3$ なら十三以下であり、重複直線の場合は追加一印を許しても十七に届かない。
$d=4$ なら $Y$ は被約でその全てが曲線部分となり、曲線外の点は二つ以下である。
従って $2d+2$（追加一印を許しても十一）に過ぎない。これで全ての被約な四次式を排除した。
\end{proof}

\begin{corollary}\label{cor:rank-two-plane-independent-curves}
同じ被約窓で、階数二平面が階数一曲線を含むことはない。
\end{corollary}
\begin{proof}
平面外印の集合はどの非零二次式の零集合にも入らない。入ればその二次式と $h$ の積が
全二十印で消え、三次評価の単射性に反する。一方、任意の非零な非被約四次式の特異集合は
ある二次式の零集合に含まれる。重複二次因子ならその二次曲面でよい。
重複平面を二回除いた残余は二次式であり、重複平面の外で特異になるのはその残余の特異部分のみである。
二次式の特異部分は一次偏導の共通零集合で、少なくとも一つの平面に入る。
$L$ の階数二以下の全点は行列式の特異点だから、前定理により $\det L=0$ が恒等的に成立する。

一般階数三なら三次余因子を $a gg^t$ と原始に書く。$g$ が定数なら固定核は全ての $L(x)e$ を消し、
この列の四独立性に反する。従って $g$ は一次列、$a$ は一次式で、階数二以下の点は
その平面と $g$ の一次基点に入る。その和も二次式で消え、不可能である。
最後に一般階数二以下なら、$h\ne0$ での消去から $N_0$ は一般階数一以下である。
その六つの二次成分は独立であるから階数一であり、原始表示は一次列 $b$ の平方の非零定数倍に限られる。
$b$ の三成分は独立である。全印で射影出力三成分と $b$ の交代積が零になり、
反対二位置の混合評価の単射性に反する。これで四行一次表示の全ての階数を尽くした。
\end{proof}
